\paragraph{Old logarithm convention.}
The affected node is \texttt{OldCoherentTreeCount}.  The prior tablet
encoding used the threshold
\[
 32tq\,\operatorname{Nat.log}(2,q),
\]
although the reference paper uses the natural logarithm in
$w=32tq\log q$.  For powers of two these quantities differ by the factor
$1/\log 2$, so the prior encoding cannot be used with the construction's
$\delta/200$ interval to prove $w\leq\delta k/5$.

The repaired Lean statement uses
\[
 32tq\,\operatorname{Nat.ceil}(\operatorname{Real.log}q).
\]
This is the natural-logarithm threshold rounded upward to an integer, which
is the appropriate exponent convention for natural $k$.  The coherent-tree
decay proof remains valid with this larger integer threshold.  In the
construction, if $x=k/s$ and
$q\leq(\delta/200)x/\log x$, then for sufficiently large $x$,
\[
 \operatorname{Nat.ceil}(\log q)
 \leq \log q+1
 \leq \log x+1
 \leq \frac54\log x.
\]
Consequently,
\[
 32tq\,\operatorname{Nat.ceil}(\log q)
 \leq \frac{\delta}{5}\frac{t}{s}k
 \leq \frac{\delta}{5}k.
\]
Thus the repaired threshold invokes the old counting estimate and returns
to the paper's stated forward-independent bound with the same
$32q^{2t-(t-1)(1-\delta/5)}$ conclusion.  No paper target or final
construction exponent is changed.
